%% Chapter 1: Introduction: chapters/chapter01.tex

% ==============================================================================
% CONTENT BLUEPRINT — Chapter 1: Introduction (M.S. Thesis)
% ==============================================================================
% Reading audience: M.S. committee, external examiner, peer researchers.
% Goal: establish the research problem and its significance in 8–12 pages.
%
% Recommended length for an M.S. thesis: 10–15 pages.
%
% ----------------------------------------------------------------------------
% §1.1 Background and Motivation  (2–3 pages)
%   • Two paragraphs to establish context:
%       – Paragraph 1: the field in general terms, why it matters now
%         (technological drivers, recent advances).
%       – Paragraph 2: the specific sub-problem the thesis addresses
%         (e.g. multi-component superlattices under biaxial strain).
%   • 6–10 references in this section, mixing classic theory papers, recent
%     reviews (last 5 years), and 2–3 closely-related experimental works.
%   • Close with a one-sentence statement of the unresolved issue, which the
%     literature review (§1.2) will substantiate.
%
% §1.2 Problem Statement  (1 paragraph)
%   • Crisp formulation: "Despite the progress summarised in §1.3, the
%     connection between [process A] and [observable B] under [conditions C]
%     has not been quantified."
%   • Optional: quantitative framing (e.g. "no published value to within
%     ±15 % at room temperature for X-class materials").
%
% §1.3 Literature Review  (3–5 pages)
%   • Organise chronologically OR thematically — pick one and state it.
%   • Cluster into 3 groups:
%       – Foundational theoretical works,
%       – Methodological developments (numerical / experimental),
%       – Material-specific studies (closest to your system).
%   • For each cluster: what was done, what was found, citation, and one-line
%     critique.
%   • Final paragraph: the gap your thesis fills, explicitly stated.
%   • Use \citep{key1,key2} for grouped citations and \citet{key} for
%     textualised citations; both work with natbib+\bibliography.
%
% §1.4 Research Objectives  (1 page)
%   • 4–6 objectives. Each bullet must be:
%       – Specific (which material, which method, which output),
%       – Testable (you can say pass/fail at the end),
%       – Bounded (achievable by thesis submission date).
%   • Bad: "To study optical properties." Good: "To compute the complex
%     dielectric function ε(ω) for monolayer MoS₂ under biaxial strain
%     from 0 % to 6 % in steps of 1 %, using PBE+SOC in Quantum ESPRESSO."
%
% §1.5 Scope and Limitations  (½–1 page)
%   • Inclusions: materials, methods, parameter ranges, code bases used.
%   • Exclusions: temperature dependence, excitonic effects beyond Tamm-
%     Dancoff, point defects, lattice vibrations, spin-orbit coupling if
%     you decide to skip it. State exclusions explicitly.
%   • Note any approximation (LDA / GGA / hybrid) and the consequence.
%
% §1.6 Thesis Organisation  (½ page)
%   • One short paragraph per subsequent chapter. State its contribution and
%     where its key results are summarised in §5.x.
%   • Forward-pointer to the §2.1 theoretical framework.
%
% Writing-style rules:
%   • Past tense for what others reported; present tense for current
%     understanding.
%   • Use \gls{key} for acronyms declared in glossary/acronyms.tex. Always
%     spell out the acronym on first use within the chapter (the glossaries
%     package handles this automatically).
%   • Equations referenced from later chapters must carry \label{eq:...}
%     and use \eqref{...}.
% ==============================================================================

\chapter{Introduction}
\label{ch:intro}

\section{Background and Motivation}
Advanced semiconductor materials play a critical role in the development of modern optoelectronic devices. In recent years, understanding the fundamental physical properties of these materials at the atomic level has become essential for developing high-performance applications. In this context, \gls{dft} offers a powerful framework to predict and tune electronic states.

\section{Problem Statement}
Despite extensive experimental studies, the electronic transitions under strain in these multi-layered systems remain poorly understood. A comprehensive first-principles study is required to model the electronic structure and clarify the exact physical mechanisms governing these systems.

\section{Research Objectives}
The primary objectives of this M.S. thesis are:
\begin{itemize}
    \item To investigate the structural stability of the chosen semiconductor materials.
    \item To calculate the electronic band structure, \gls{dos}, and \gls{pdos}.
    \item To model the effects of biaxial strain on the bandgap and carrier effective masses.
    \item To compute the optical properties, including the complex dielectric function and absorption coefficient.
\end{itemize}

\section{Scope of the Study}
This study is focused on the ground-state properties of bulk and monolayer configurations using the \gls{vasp} and Quantum ESPRESSO code bases. Temperature-dependent properties and excitonic effects are beyond the scope of the present work.

\section{Thesis Organization}
This thesis is organized into six chapters. Chapter~\ref{ch:intro} introduces the background, objective, and scope. Chapter~\ref{ch:lit_review} provides a review of prior work. Chapter~\ref{ch:methodology} describes the theoretical framework and computational setup. Chapter~\ref{ch:results_1} presents the structural and electronic results. Chapter~\ref{ch:results_2} details the optical properties and discussion. Finally, Chapter~\ref{ch:conclusion} summarizes findings and outlines future directions.
